% ---------------------------------------------------------------------------
% 1 INTRODUÇÃO
% ---------------------------------------------------------------------------
\chapter{Introdução}\label{cap:introducao}

A Introdução é a apresentação do tema a ser tratado e deve conter o problema
a ser pesquisado.

Ao desenvolver a introdução, o grupo deve explicar o assunto que deseja
abordar, de forma a:

\begin{itemize}
    \item Desenvolver o tema;
    \item Anunciar a ideia básica;
    \item Delimitar o foco da pesquisa;
    \item Situar o tema dentro do contexto geral da sua área de trabalho;
    \item Descrever as motivações que levaram à escolha do tema;
    \item Indicar o objeto do trabalho: o que será estudado?
\end{itemize}

O texto do trabalho deve conter a formatação indicada neste documento:

\begin{itemize}
    \item \textbf{Fonte, tamanho e cor:} Times New Roman, tamanho 12 para
          texto, 10 para citações de mais de três linhas e de 10 para notas
          de rodapé; Cor preta.
    \item \textbf{Margens:} superior e esquerda de 3cm; inferior e direita
          de 2cm.
    \item \textbf{Títulos ou subtítulos:} alinhados à esquerda, iniciando
          sempre em uma nova página. Todas as letras dos títulos dos
          capítulos devem ser escritas no canto esquerdo de cada página, em
          negrito e em maiúsculas.
    \item \textbf{Paginação} (números das páginas): superior à direita,
          começando na introdução em algarismos arábicos (1, 2, 3\ldots).
    \item \textbf{Espaçamento:} Todo o texto deve ser digitado em espaço
          1,5. Excetuam-se: citações longas (com mais de três linhas), notas
          de rodapé, as Referências Bibliográficas (ou Bibliografia) e as
          legendas de ilustrações e tabelas, que são digitadas com
          espaçamento simples. Os parágrafos devem ser separados por uma
          linha em branco. Citações com mais de três linhas: fonte tamanho
          10, espaçamento simples e recuo de 4cm da margem esquerda. Notas
          de rodapé: fonte tamanho 10.
\end{itemize}
